\chapter{Das Raytracer-Plugin (Raytatouille)}\label{kap:raytracer}
\index{Raytracer}\index{Raytatouille}

Das Plugin \plugin{Raytracer} verbindet \produkt{} mit dem optischen Rechenkern
\emph{Raytatouille}. Raytatouille beschreibt ein optisches System vollständig in einer Datei
(\datei{.rtt.json}): Flächen mit ihren exakten Formen, Gläser, Blenden, Felder,
Wellenlängen, Strahlengänge (\emph{Paths}) und auf Wunsch eine Optimierungsaufgabe. Das Plugin
lädt ein solches System, zeichnet jede optische Fläche als NURBS-Fläche in das CAD-Modell,
verfolgt Strahlen, berechnet die Bestrahlungsstärke auf der Bildfläche, optimiert das System mit
Live-Aktualisierung der Geometrie und schreibt das Ergebnis wieder als \datei{.rtt.json}.

Dieses Kapitel beschreibt die Arbeit mit dem Plugin aus Anwendersicht. Die vollständige
Referenz aller Aktionen, Parameter, Ergebnisse und Attribute finden Sie in
\autoref{kap:plugin-raytracer}. Wie Plugins allgemein eingerichtet und bedient werden, steht in
\autoref{kap:plugins}; der Node-Workflow, in dem das Plugin Blöcke anbietet, in
\autoref{kap:workflow}.

\abbildung[0.95\linewidth]{raytracer.png}{Das Raytracer-Plugin: ein Singlet mit
Strahlenbündel in allen vier Ansichten, rechts das Panel \gui{Plugin} mit den Parametern der
Gruppe \gui{System} und den Aktionsknöpfen, unten in der Statuszeile die Meldung der letzten
Aktion.}{fig:raytracer-uebersicht}

\section{Grundidee: Das optische System ist maßgeblich}\label{sec:rt-grundidee}

Beim Raytracer gilt ein einfacher Grundsatz: \textbf{Maßgeblich ist das optische System, nicht
die Zeichnung.} Die NURBS-Flächen im CAD-Modell werden aus den exakten Flächen des Systems
abgeleitet -- nie umgekehrt. Strahlen werden immer auf den exakten Flächen von Raytatouille
verfolgt, nicht auf den gezeichneten NURBS. Die Zeichnung dient zur Anschauung, zur Kontrolle der
Lage im Raum und als Träger der Ergebnisse (Strahlen, Bestrahlungsstärke als Farbe auf der
Bildfläche).

Daraus folgt:
\begin{itemize}
  \item Wenn Sie eine gezeichnete Linsenfläche mit den CAD-Werkzeugen verschieben oder
    verformen, ändert das die Rechnung nicht. Optische Daten (Radien, Dicken, Gläser) ändern Sie
    in der \datei{.rtt.json}-Datei oder -- wenn ein Node-Workflow das System enthält -- im Dialog
    \gui{Parameters…} der Blöcke (\autoref{sec:rt-workflow}).
  \item Jede Aktion, die das System ändert (Laden, Optimieren), zeichnet die betroffenen
    Flächen neu.
  \item Es gibt zwei mögliche Quellen des Systems: die geladene Datei (\gui{System file}) oder
    den Node-Workflow des Dokuments. Sobald ein Workflow das System enthält, ist \emph{er}
    maßgeblich (\autoref{sec:rt-workflow}).
\end{itemize}

\section{Voraussetzungen}\label{sec:rt-voraussetzungen}

Das Plugin wird mit der Anwendung ausgeliefert. Es braucht im Python-Interpreter der Plugins
(\menu{Plugins,Settings…}) folgende Pakete; die Namen stehen so in der Datei
\datei{plugin.json} des Plugins:

\begin{referenz}{@{}l l >{\raggedright\arraybackslash}X@{}}
\toprule
Import-Name & pip-Paket & Wofür \\
\midrule\endhead
\code{numpy} & \code{numpy} & Zahlenfelder (immer nötig) \\
\code{raytatouille} ($\geq$ 0.9) & \code{raytatouille} & der optische Rechenkern: Laden, Strahlen, Bestrahlungsstärke, Optimierung \\
\code{rtt\_fusion} ($\geq$ 0.2) & \code{rtt-fusion} & Node-Workflow: \gui{New Raytatouille system}, \gui{Workflow from system}, Blöcke \code{rtt.*} \\
\code{fieldtracing} ($\geq$ 0.2) & \code{photonpower-fieldtracing} & Felder (Darstellung H) in Verbindungen des Workflows \\
\bottomrule
\end{referenz}

\gui{Check environment} in \menu{Plugins,Settings…} zeigt, was fehlt, mit dem passenden
\code{pip install}-Befehl; \gui{Create environment…} legt eine neue Umgebung an und installiert
numpy und die Pakete aller geladenen Plugins (siehe \autoref{kap:installation} und
\autoref{kap:plugins}).

\begin{hinweis}
Fehlt \code{raytatouille}, startet das Plugin trotzdem; erst eine Aktion meldet dann:
\enquote{\texttt{The Python package 'raytatouille' is not installed in the plugin's Python
environment. Plugins $\to$ Settings… $\to$ Check environment shows how to install it.}}
Fehlt nur \code{rtt\_fusion}, funktionieren Laden, Strahlen, Kamerabild, Optimieren und
Speichern; nur die Workflow-Aktionen melden einen Fehler. Ohne \code{fieldtracing} kann der
Workflow nur mit Strahlen (Darstellung R) rechnen.
\end{hinweis}

Für die Bestrahlungsstärke braucht das Plugin eine Raytatouille-Version mit der Funktion
\code{rt.analysis.irradiance\_map}. Bei einer zu alten Version meldet es: \enquote{\texttt{The
irradiance needs a newer Raytatouille (rt.analysis.irradiance\_map, \#379).}}

\section{Das Raytracer-Panel}\label{sec:rt-panel}
\index{Raytracer!Panel}

Wählen Sie im Dock \gui{Plugin} oben in der Auswahlliste den Eintrag \gui{Raytracer}. Das Panel
zeigt die Beschreibung des Plugins, darunter die Parameter in vier Gruppen und darunter die
Aktionsknöpfe. Dieselben Aktionen stehen im Menü \menu{Plugins,Raytracer} und als Knöpfe in der
Gruppe \gui{Raytracer} der Werkzeugspalte (dort mit den ersten drei Buchstaben, z.\,B.
\gui{Loa} für \gui{Load system}).

\begin{description}
  \item[\gui{System}] Systemdatei, Glaskataloge, Zeichentoleranz.
  \item[\gui{Trace}] Auswahl der Strahlengänge, Felder und Wellenlängen, Zahl der gezeichneten
    Strahlen, Einfärbung, schnelle Bestrahlungsstärke.
  \item[\gui{Camera image}] Strahlenzahl und Pixel des Kamerabilds.
  \item[\gui{Optimise}] Zieldatei für \gui{Save system}.
\end{description}

Die Aktionen sind:

\begin{referenz}{@{}l >{\raggedright\arraybackslash}X@{}}
\toprule
Aktion & Kurzbeschreibung \\
\midrule\endhead
\gui{Load system} & liest die \datei{.rtt.json} und zeichnet sie (ersetzt eine früher geladene Zeichnung) \\
\gui{Trace rays} & Strahlen entlang der Strahlengänge als Linien; auf Wunsch schnelle Bestrahlungsstärke \\
\gui{Camera image} & Bestrahlungsstärke auf der letzten Fläche jedes Strahlengangs mit vielen Strahlen als Bild \\
\gui{Optimise} & Optimierer von Raytatouille; die Flächen werden nach jeder Iteration neu gezeichnet \\
\gui{Save system} & schreibt das System als \datei{.rtt.json} \\
\gui{New Raytatouille system} & neues System ohne Datei im Node-Workflow: Quelle, Blende, Detektor \\
\gui{Workflow from system} & Node-Workflow aus dem geladenen System \\
\gui{Update drawing} & zeichnet das System des Node-Workflows jetzt \\
\gui{Follow the workflow} & zeichnet, was eine Änderung im Workflow geändert hat (läuft automatisch) \\
\gui{Section $\to$ group} & nur im Kontextmenü eines Blocks \code{rtt.section} im Node-Workflow \\
\bottomrule
\end{referenz}

Die Parameterwerte werden \textbf{mit dem Dokument gespeichert} (pro Plugin). Eine Änderung
eines Parameters markiert das Dokument als geändert, ist aber kein Rückgängig-Schritt.
\gui{Reset parameters} unten im Panel setzt alle Parameter auf ihre Standardwerte zurück -- also
auch die Systemdatei und die Liste der Glaskataloge.

Jede Aktion läuft im Hintergrund; der Fortschritt steht in der Statuszeile, und \keys{Esc} oder
\gui{Cancel} bricht ab. Nach dem Ende steht eine kurze Zusammenfassung in der Statuszeile;
ausführlichere Hinweise (z.\,B. welche Glaskataloge geladen wurden) schreibt das Plugin in das
Dock \gui{Plugin log}. Alle Änderungen einer Aktion am Modell bilden \textbf{einen}
Rückgängig-Schritt.

\section{Ein optisches System laden}\label{sec:rt-laden}
\index{Raytracer!System laden}\index{rtt.json@\texttt{.rtt.json}}

\subsection{Die Systemdatei}

Ein Raytatouille-System ist eine JSON-Datei mit der Endung \datei{.rtt.json}. Das Plugin
bringt ein kleines Beispiel mit, \datei{examples/singlet\_optim.rtt.json} im Ordner des Plugins:
eine Einzellinse aus einem Glas mit konstantem Brechungsindex 1{,}5168, die auf eine Brennweite
von \SI{100}{\milli\metre} und minimalen RMS-Fleck optimiert werden soll. An ihm lässt sich der
Aufbau einer Systemdatei gut ablesen:

\begin{lstlisting}
{
  "schema_version": "0.9.0",
  "name": "Singlet to optimize: EFL 100 mm and minimal RMS spot, ...",
  "units": {"length": "mm", "wavelength": "um"},
  "wavelengths": [ {"um": 0.5876, "reference": true} ],
  "aperture": {"type": "epd", "value": 10.0},
  "fields": { "points": [ {} ] },
  "root": {
    "type": "assembly", "name": "system",
    "children": [
      { "type": "stop", "name": "stop",
        "surfaces": [ {"id": "STO",
                       "aperture": {"type": "circular", "radius": 5.0}} ] },
      { "type": "lens", "name": "L1",
        "pose": {"position": [0.0, 0.0, 5.0]},
        "material": "CONST:1.5168",
        "surfaces": [ ... "radius": {"value": 103.36, "variable": true} ... ] },
      { "type": "detector", "name": "image", ... "surfaces": [ {"id": "IMG"} ] }
    ]
  },
  "paths": [ {"name": "main", "events": "auto"} ],
  "optimization": {
    "operands":   [ {"type": "efl", "path": "main", "target": 100.0,
                     "weight": 10000.0} ],
    "generators": [ {"type": "rms_spot", "path": "main"} ]
  }
}
\end{lstlisting}

Die wichtigsten Teile:
\begin{itemize}
  \item \param{wavelengths}: die Wellenlängen in \si{\micro\metre}, eine davon die
    Referenzwellenlänge.
  \item \param{aperture}: die Apertur des Systems (hier: Eintrittspupillendurchmesser
    \SI{10}{\milli\metre}).
  \item \param{fields}: die Feldpunkte (hier einer auf der Achse).
  \item \param{root}: der Baum aus Baugruppen (\code{assembly}) und Elementen (Blende, Linse,
    Detektor, \dots) mit ihren Flächen. Jede Fläche hat eine eindeutige \param{id}
    (\code{STO}, \code{L1.S1}, \dots).
  \item \param{paths}: die Strahlengänge mit Namen (hier \code{main}). Ein System kann mehrere
    haben, z.\,B. die beiden Arme eines Interferometers.
  \item \param{optimization}: was der Optimierer anstrebt; Werte mit
    \code{\dq{}variable\dq{}: true} sind Variablen der Optimierung.
\end{itemize}
Das Format selbst gehört zu Raytatouille und ist dort beschrieben; dieses Handbuch beschreibt
nur, wie \produkt{} damit umgeht.

\subsection{Schritt für Schritt}

\begin{schritte}
  \item Öffnen Sie das Dock \gui{Plugin} und wählen Sie \gui{Raytracer}.
  \item Klicken Sie in der Gruppe \gui{System} neben \gui{System file} auf den Knopf zur
    Dateiauswahl und wählen Sie die \datei{.rtt.json}-Datei. Der Dateidialog zeigt
    \enquote{Raytatouille system (*.rtt.json *.json)}.
  \item Wenn das System Gläser aus Katalogen verwendet: Fügen Sie die Kataloge unter
    \gui{Glass catalogues} hinzu oder legen Sie sie neben die Systemdatei
    (\autoref{sec:rt-glaskataloge}).
  \item Klicken Sie auf \gui{Load system}.
\end{schritte}

Das Plugin lädt die Glaskataloge, liest und übersetzt (\enquote{kompiliert}) das System und
zeichnet jede Fläche. In der Statuszeile erscheint der Name des Systems und die Zahl der Flächen,
z.\,B. \enquote{\texttt{Singlet to optimize: … : 4 surfaces}}. Enthält das Dokument noch keine
Geometrie, zoomt die Anwendung alle Ansichten auf das Ergebnis.

\gui{Load system} \textbf{ersetzt} eine früher geladene Zeichnung: Alle Objekte, die das Plugin
gezeichnet hat (erkennbar an den Attributen \param{rtt.surface} oder \param{rtt.rim}), und alle
Gruppen, deren Name mit \enquote{\texttt{Raytatouille: }} beginnt, werden zuerst entfernt.
Andere Objekte des Dokuments -- etwa eine Fassung, die Sie selbst konstruiert haben -- bleiben
unberührt.

\subsection{Was gezeichnet wird}

\begin{itemize}
  \item \textbf{Jede optische Fläche} wird eine NURBS-Fläche, benannt nach ihrer \param{id}
    (z.\,B. \code{L1.S1}). Die Begrenzung (Apertur) ist eine exakte Trimmkurve: Kreis oder Ellipse,
    gegebenenfalls mit Loch in der Mitte; Rechtecke brauchen keine Trimmkurve.
  \item \textbf{Linsenränder}: Bei Linsen mit zwei Flächen und kreisförmigen Aperturen gleichen
    Radius wird der Rand als Regelfläche zwischen den beiden Randkurven gezeichnet. Diese
    Objekte heißen \enquote{\textit{Element} edge} (bei mehreren \enquote{edge 2} usw.).
  \item \textbf{Die Aperturblende} erscheint als Ring.
  \item \textbf{Die Bildfläche} (Detektor) hat in vielen Systemen keine eigene Apertur. Dann
    bestimmt das Plugin ihre Größe aus einer kurzen Probe-Strahlverfolgung: aus der Ausdehnung des
    Flecks plus \SI{25}{\percent} Rand, mindestens aber \SI{5}{\percent} des größten
    Aperturradius des Systems -- ein scharfer Fokus ergäbe sonst eine Fläche, die neben dem
    System nicht zu sehen wäre.
\end{itemize}

\subsection{Zeichengenauigkeit}\label{sec:rt-toleranz}
\index{Zeichentoleranz}

Der Parameter \gui{Drawing tolerance} (Standard \SI{0.0001}{\milli\metre} =
\SI{0.1}{\micro\metre}) legt fest, wie weit die gezeichnete NURBS-Fläche höchstens von der
optischen Fläche abweichen darf. Er betrifft \emph{nur die Zeichnung}; Strahlen laufen immer
auf den exakten Flächen.

\begin{itemize}
  \item \textbf{Rotationssymmetrische Flächen} (Kegelschnitte, Asphären, Axicon, Q-Typ) mit
    kreisförmiger Apertur werden als Rotationsflächen gezeichnet: Das Profil wird entlang
    seiner Bogenlänge interpoliert, sodass auch steile, fast halbkugelförmige Linsen so genau
    gezeichnet werden wie flache.
  \item \textbf{Andere Formen}, rechteckige oder elliptische Aperturen und Bildflächen werden
    als bikubische B-Spline-Flächen über der lokalen x-y-Ebene gezeichnet; das Kontrollnetz wird
    verfeinert, bis die Toleranz eingehalten ist (6, 10, 18, 34 oder 66 Kontrollpunkte je
    Richtung).
  \item \textbf{Ebenen} sind exakt.
\end{itemize}

Die tatsächliche Abweichung jeder Fläche steht in ihrem Attribut \param{rtt.nurbs\_error\_mm}
(sichtbar im Dock \gui{Properties}, wenn Sie die Fläche auswählen). Flächen, die die Toleranz
überschreiten -- typisch bei nicht glatten Formen wie der Spitze eines Axicons --, nennt das
Plugin im \gui{Plugin log}: \enquote{\texttt{drawn with more than the tolerance (non-smooth
shapes such as the axicon tip): …}}.

\begin{tipp}
Eine größere Toleranz (z.\,B. \SI{0.001}{\milli\metre}) macht das Zeichnen großer Systeme
schneller und die Flächen leichter, ohne die Rechnung zu verändern. Eine kleinere lohnt sich
nur, wenn Sie die gezeichneten Flächen weiterverwenden, z.\,B. als STEP für die Konstruktion einer
Fassung exportieren.
\end{tipp}

\section{Der Modellbaum des Systems}\label{sec:rt-modellbaum}
\index{Modellbaum!Raytracer}

Das Plugin legt im Dock \gui{Model} eine Gruppenstruktur an, die dem Baum des Systems folgt:

\begin{itemize}
  \item eine oberste Gruppe \enquote{\texttt{Raytatouille: }\textit{Name}} -- \textit{Name} ist
    der Name des Systems aus der Datei oder, wenn dieser fehlt, der Dateiname ohne Endung; ein
    aus dem Workflow gezeichnetes System heißt \enquote{\textit{Name} (workflow)};
  \item darunter je eine Gruppe pro Baugruppe (\code{assembly}) und pro Element (Linse, Blende,
    Spiegel, Detektor, \dots), benannt nach dem \param{name} des Knotens, sonst nach seinem Typ;
  \item in der Gruppe eines Elements dessen Flächen und Linsenränder.
\end{itemize}

So können Sie mit den Kontrollkästchen des Modellbaums einzelne Linsen oder ganze Baugruppen
ausblenden, etwa um in ein Objektiv hineinzusehen. Ausblenden ändert die Rechnung nicht.

Jede gezeichnete Fläche trägt Attribute, die Sie im Dock \gui{Properties} sehen:

\begin{referenz}{@{}l >{\raggedright\arraybackslash}X@{}}
\toprule
Attribut & Bedeutung \\
\midrule\endhead
\param{rtt.surface} & \param{id} der Fläche im System \\
\param{rtt.location} & Ort der Fläche im Systembaum (JSON-Pfad, z.\,B. \code{/root/children/1/surfaces/0}) \\
\param{rtt.kind} & Art der Fläche laut Raytatouille (z.\,B. Blende, Detektor) \\
\param{rtt.shape} & Form der Fläche laut Raytatouille \\
\param{rtt.nurbs\_error\_mm} & größte Abweichung der Zeichnung in mm \\
\param{rtt.anchor} & nur an der ersten Fläche: Anker des Systems \\
\param{rtt.revision} & nur am Anker: Nummer des Systemstands (\autoref{sec:rt-undo}) \\
\param{rtt.file} & nur am Anker: Pfad der geladenen Datei \\
\param{rtt.workflow} & nur am Anker, wenn das System aus dem Workflow stammt: Kennung seines Stands \\
\param{rtt.rim} & an Linsenrändern: Name des Elements \\
\bottomrule
\end{referenz}

\begin{warnung}
Ändern oder löschen Sie diese Attribute nicht von Hand. Das Plugin erkennt seine Zeichnung und
den aktuellen Systemstand an ihnen. Fehlt der Anker, lädt das Plugin bei der nächsten Aktion die
Systemdatei neu.
\end{warnung}

\section{Glaskataloge}\label{sec:rt-glaskataloge}
\index{Glaskatalog}\index{AGF}

Gläser wie \code{N-BK7} oder \code{N-SF6} kommen aus Glaskatalogen im AGF-Format
(\datei{*.agf}). Ein Material mit konstantem Brechungsindex (\code{CONST:1.5168} im Beispiel)
braucht keinen Katalog.

\subsection{Woher das Plugin die Kataloge nimmt}

Bei jedem Laden verwendet das Plugin, in dieser Reihenfolge und ohne Doppelte:
\begin{enumerate}
  \item die Dateien der Liste \gui{Glass catalogues} im Panel (Gruppe \gui{System});
  \item alle Dateien \datei{*.agf} im Ordner der Systemdatei;
  \item alle Dateien \datei{*.agf} im Unterordner \datei{catalogs} neben der Systemdatei;
  \item Kataloge, die im Node-Workflow des Dokuments gespeichert sind, soweit keine Datei sie
    liefert (siehe unten).
\end{enumerate}

Nach dem Laden nennt das \gui{Plugin log} alle verwendeten Kataloge in einer Zeile wie
\enquote{\texttt{glass catalogues: SCHOTT.agf, OHARA.agf}}; gespeicherte Kataloge tragen den
Zusatz \enquote{\texttt{(stored)}}. Ein Katalog, der sich nicht lesen lässt, wird mit seiner
Fehlermeldung im Log genannt (\enquote{\texttt{catalogue …: …}}) und übersprungen; die anderen
werden trotzdem geladen.

\subsection{Kataloge hinzufügen und entfernen}

\begin{schritte}
  \item Klicken Sie unter \gui{Glass catalogues} auf \gui{Add…}.
  \item Wählen Sie eine oder mehrere AGF-Dateien (der Filter zeigt \enquote{Glass catalogues
    (*.agf *.AGF)}, wahlweise \enquote{All files (*)}).
  \item Zum Entfernen markieren Sie einen Eintrag und klicken auf \gui{Remove}.
  \item Laden Sie das System neu (\gui{Load system}), damit die geänderte Liste gilt.
\end{schritte}

\subsection{Kataloge im Dokument}

Die Liste \gui{Glass catalogues} ist ein Parameter des Panels und wird daher mit dem Dokument
gespeichert -- allerdings nur als Liste von \emph{Pfaden}. Damit ein Dokument auch auf einem
anderen Rechner ohne die AGF-Dateien rechnet, speichert der Node-Workflow die Kataloge selbst:
Das Panel des Raytracers ist die \emph{eine} Stelle, an der Sie Kataloge angeben. Der Workflow
übernimmt die dort gelisteten Kataloge (und die \datei{*.agf} neben der Systemdatei) und legt
ihren Inhalt bei der nächsten Prüfung oder Rechnung (\gui{Check workflow}, \gui{Run workflow})
im Quellblock des Systems ab. Danach rechnet die \datei{.rttf}-Datei ohne die AGF-Dateien, und
auch das Raytracer-Panel verwendet die gespeicherten Kataloge, wenn eine Datei fehlt.

\begin{tipp}
Wenn Sie ein Dokument weitergeben wollen, das Gläser aus Katalogen verwendet: Legen Sie mit
\gui{Workflow from system} einen Workflow an, führen Sie einmal \gui{Check workflow} aus und
speichern Sie das Dokument. Die Kataloge reisen dann im Dokument mit.
\end{tipp}

\subsection{Typische Fehler mit Katalogen}

Fehlt ein Glas, bricht \gui{Load system} mit einer Meldung der Form \enquote{\textit{Datei}
\texttt{cannot be used:} \textit{Fehler von Raytatouille}} ab. Erwähnt die Fehlermeldung einen
Katalog, ergänzt das Plugin den Hinweis \enquote{\texttt{Glass catalogues: put the *.agf files
next to the system file or list them under Glass catalogues.}} Prüfen Sie dann:
\begin{itemize}
  \item ob der Katalog, aus dem das Glas stammt, in der Liste oder neben der Systemdatei liegt;
  \item ob das Log den Katalog als geladen nennt -- wenn nicht, steht dort der Grund;
  \item ob der Glasname in der Systemdatei genau so im Katalog steht.
\end{itemize}

\section{Quelle: Felder, Wellenlängen und Apertur}\label{sec:rt-quelle}

Was in das System hineinleuchtet, legt die Systemdatei fest: die Feldpunkte (\param{fields}),
die Wellenlängen (\param{wavelengths}) mit ihrer Referenzwellenlänge und ihren Gewichten, die
Apertur (\param{aperture}) und die Lage des Objekts (im Unendlichen oder in endlichem Abstand).
Im Node-Workflow erscheinen diese Einstellungen im Quellblock \block{rtt.source}; dessen
\gui{Parameters…} zeigt sie zum Bearbeiten (\autoref{sec:rt-workflow}).

Im Panel wählen Sie für \gui{Trace rays} und \gui{Camera image} nur aus, \emph{welche} Felder und
Wellenlängen gerechnet werden:

\begin{description}
  \item[\gui{Field}] \code{0} rechnet alle Feldpunkte; \code{1}, \code{2}, \dots{} nur den
    ersten, zweiten, \dots{} Feldpunkt des Systems.
  \item[\gui{Wavelengths}] \code{all} rechnet alle Wellenlängen des Systems, \code{reference}
    nur die Referenzwellenlänge.
\end{description}

Die Strahlen starten auf einem hexapolaren Raster über der Pupille: ein Strahl in der Mitte und
Ringe mit 6, 12, 18, \dots{} Strahlen. Deshalb sind nur bestimmte Strahlenzahlen möglich --
1, 7, 19, 37, 61, 91, 127, \dots{} (allgemein $1 + 3r(r+1)$ bei $r$ Ringen). Das Plugin nimmt
die Ringzahl, deren Strahlenzahl der eingestellten am nächsten liegt, mindestens aber einen
Ring.

\section{Strahlen verfolgen}\label{sec:rt-strahlen}
\index{Raytracer!Strahlen verfolgen}\index{Strahlengang}

\gui{Trace rays} verfolgt Strahlen entlang der Strahlengänge des Systems und zeichnet sie als
Linien in die Ansichten.

\begin{schritte}
  \item Laden Sie ein System (\autoref{sec:rt-laden}).
  \item Stellen Sie in der Gruppe \gui{Trace} ein:
    \begin{itemize}
      \item \gui{Paths}: die Namen der Strahlengänge, durch Kommas getrennt, z.\,B.
        \code{reference, sample}. Leer: alle Strahlengänge des Systems.
      \item \gui{Field} und \gui{Wavelengths} wie in \autoref{sec:rt-quelle}.
      \item \gui{Rays (drawn)}: Strahlen je Feld und Wellenlänge (Standard 127, höchstens
        20\,000).
      \item \gui{Colour by}: Einfärbung nach \code{wavelength}, \code{field} oder \code{path}.
    \end{itemize}
  \item Klicken Sie auf \gui{Trace rays}.
\end{schritte}

\subsection{Ergebnis}

Im Dock \gui{Results} entsteht ein neuer Lauf \enquote{\texttt{Trace }\textit{Datei}}
(\textit{Datei} ist der Name der Systemdatei bis zum ersten Punkt, z.\,B. \code{Trace
singlet\_optim}). Er enthält für jeden Strahlengang einen Linien-Satz \enquote{\texttt{Rays
}\textit{Pfad}}, den Sie einzeln ein- und ausblenden können. Jeder Strahl ist eine Polylinie von
Fläche zu Fläche; ein Strahl, der unterwegs verloren geht (vignettiert, Totalreflexion, \dots),
endet dort.

Die Einfärbung:
\begin{referenz}{@{}l l >{\raggedright\arraybackslash}X@{}}
\toprule
\gui{Colour by} & Farbskala & Wert \\
\midrule\endhead
\code{wavelength} & viridis & Wellenlänge in nm \\
\code{field} & plasma & Nummer des Feldpunkts (1, 2, \dots) \\
\code{path} & coolwarm & Nummer des Strahlengangs (1 bis Zahl der Strahlengänge) \\
\bottomrule
\end{referenz}

Die Statuszeile meldet für jeden Strahlengang, wie viele der gezeichneten Strahlen sein Ende
erreichen, z.\,B. \enquote{\texttt{rays reaching the end of the path – main: 217/217}}. Weniger
als alle bedeutet Vignettierung oder andere Verluste.

\begin{tipp}
Einzelne Strahlen lassen sich in CAD-Kurven umwandeln, etwa um daran zu messen oder sie zu
exportieren: Rechtsklick auf den Linien-Satz im Dock \gui{Results}, \gui{Convert to CAD
curves…}, dann im Dialog \enquote{Every n-th ray} angeben, jeder wievielte Strahl übernommen
wird.
\end{tipp}

\subsection{Strahlen anklicken, verschieben, drehen, löschen}\label{sec:rt-strahlen-objekt}

Die berechneten Strahlen eines Laufs sind in den Ansichten \emph{ein} Element: Ein Klick auf
einen beliebigen Strahl wählt alle Strahlen dieses Laufs aus (gelb hervorgehoben, unter dem
Mauszeiger türkis). \gui{Move} und \gui{Rotate} verschieben bzw.\ drehen sie, \keys{Entf}
löscht sie -- jeweils rückgängig machbar (Abschnitt~\ref{sec:ergebnis-zeichnungen}).

\begin{hinweis}
Verschieben und Drehen von Strahlen ist eine reine \emph{Darstellungs}-Transformation: Die
Strahlen werden nicht neu berechnet, und das optische System ändert sich nicht. So lassen sich
z.\,B.\ zwei Läufe nebeneinander vergleichen. Gespeichert (\gui{Save…}) und in CAD-Kurven
umgewandelt werden sie dort, wo sie gezeigt werden.
\end{hinweis}

\subsection{Mehrere Strahlengänge}

Systeme mit Strahlteilern -- etwa ein Interferometer -- haben mehrere Strahlengänge. Ohne Angabe
unter \gui{Paths} verfolgt \gui{Trace rays} alle; jeder Strahlengang bekommt seinen eigenen
Linien-Satz. Ein unbekannter Name führt zu der Meldung \enquote{\texttt{Unknown path …. Paths of
this system: …}}, die die vorhandenen Namen aufzählt. Scheitert bei mehreren Strahlengängen einer,
werden die übrigen trotzdem verfolgt; die Statuszeile nennt den fehlgeschlagenen
(\enquote{\texttt{not traced: …}}), das Log den Grund.

\subsection{Gefaltete und gekippte Strahlengänge}\label{sec:rt-gefaltet}
\index{Strahlengang!gefaltet}

Das Zielen der Strahlen in Raytatouille (\emph{ray aiming}) braucht einen rotationssymmetrischen
Strahlengang. Für gefaltete oder gekippte Strahlengänge (Strahlteiler, Interferometer wie ein
Linnik) startet das Plugin die Strahlen selbst: auf demselben hexapolaren Raster über der
Blendenöffnung, parallel entlang der Feldwinkel (Objekt im Unendlichen) oder von den
Objektpunkten aus (Objekt in endlichem Abstand). Das ist exakt, wenn die Aperturblende die
\emph{erste} Fläche des Strahlengangs ist, wie bei einer Beleuchtungsblende. Sonst meldet das
Plugin: \enquote{\texttt{Raytatouille cannot aim rays into this path (…). The plugin can start
them itself only when the aperture stop is the first surface of the path.}}

Diese eigenen Startstrahlen gibt es nur für Feldwinkel bei unendlichem Objekt und für
Objekthöhen bei endlichem Objekt; für andere Feldtypen erscheint eine entsprechende Meldung.

\section{Bestrahlungsstärke auf der Bildfläche}\label{sec:rt-bestrahlung}
\index{Bestrahlungsstärke!Raytracer}\index{Kamerabild!Raytracer}

Das Plugin berechnet die Bestrahlungsstärke auf der \emph{letzten Fläche} jedes Strahlengangs,
in der Regel der Bildfläche. Die Rechnung übernimmt Raytatouille: Jede Zelle eines Rasters
erhält das auftreffende Strahlgewicht pro projizierter Zellfläche. Das Ergebnis ist relativ
(Einheit \code{W/mm\textsuperscript{2} (rel.)}), die Gewichte der Wellenlängen aus dem System
werden berücksichtigt. Strahlengänge, die auf derselben Fläche enden, werden addiert --
inkohärent, also ohne Interferenz.

Es gibt zwei Wege:

\subsection{Schnelle Bestrahlungsstärke mit \gui{Trace rays}}

Ist in der Gruppe \gui{Trace} \gui{Quick irradiance} eingeschaltet (Standard: aus), berechnet
\gui{Trace rays} zusätzlich eine grobe Bestrahlungsstärke mit wenigen Strahlen:
\gui{Rays (quick irradiance)} (Standard 5\,000 je Feld und Wellenlänge) auf einem Raster von
\gui{Quick irradiance grid} Zellen an der langen Seite (Standard 128). Das Ergebnis erscheint als
Farbe auf der Fläche (\enquote{\texttt{Irradiance }\textit{Fläche}} im Lauf), mit Farbleiste in
den Ansichten. Die Statuszeile nennt den Spitzenwert, z.\,B. \enquote{\texttt{irradiance on IMG
(peak 1.5e+06)}}.

\subsection{Kamerabild mit \gui{Camera image}}

Für ein glattes, aussagekräftiges Bild verwenden Sie \gui{Camera image}. Es verfolgt viele
Strahlen und zeigt das Ergebnis als Bild im Fenster \gui{Image}.

\begin{schritte}
  \item Wählen Sie in der Gruppe \gui{Trace} die Strahlengänge, Felder und Wellenlängen
    (dieselben Parameter wie für \gui{Trace rays}).
  \item Stellen Sie in der Gruppe \gui{Camera image} ein:
    \begin{itemize}
      \item \gui{Rays per field and wavelength} (Standard 200\,000, von 1\,000 bis
        50\,000\,000). Die Rechenzeit wächst linear mit der Strahlenzahl.
      \item \gui{Pixels (long side)} (Standard 512). Die Pixel sind quadratisch; der Sensor ist
        die Apertur der Fläche (bzw. die vom Plugin bestimmte Größe der Bildfläche).
      \item \gui{Also on the surface in 3D} (Standard: an): das Bild zusätzlich als Farbe auf der
        Fläche in den Ansichten.
    \end{itemize}
  \item Klicken Sie auf \gui{Camera image}.
\end{schritte}

Das Fenster \gui{Image} zeigt \enquote{\texttt{Camera image }\textit{Fläche}} mit Achsen in mm,
Farbleiste und dem Wert unter dem Mauszeiger. \gui{log} schaltet auf eine logarithmische Skala
um -- nützlich, um schwaches Streulicht neben einem hellen Fleck zu sehen. Das Mausrad zoomt,
Ziehen verschiebt, ein Doppelklick zeigt das ganze Bild. Die Beschreibung des Bildes nennt die
Fläche, die Zahl der Strahlen, die beteiligten Strahlengänge, Felder und Wellenlängen, die
Pixelzahl und die Pixelgröße in \si{\micro\metre}.

\begin{hinweis}
Strahlen, die nicht gezählt werden -- außerhalb des Rasters, vignettiert oder nicht bis zur
Fläche gelangt --, listet das \gui{Plugin log} je Fläche und Strahlengang auf
(\enquote{\texttt{rays not counted – outside the grid: …, vignetted: …, not reached: …}}).
Viele Strahlen \enquote{outside the grid} bedeuten meist, dass die Bildfläche kleiner ist als
der Fleck.
\end{hinweis}

\subsection{Einschränkungen}

\begin{itemize}
  \item Die Bestrahlungsstärke wird nur auf \textbf{ebenen} letzten Flächen abgebildet. Auf einer
    gekrümmten Rotationsfläche meldet das Log \enquote{\texttt{irradiance on …: curved surface of
    revolution – not mapped yet}}. Endet keiner der Strahlengänge auf einer gezeichneten ebenen
    Fläche, bricht \gui{Camera image} mit \enquote{\texttt{No camera image: none of the paths
    ends on a drawn plane surface.}} ab.
  \item Ändert das letzte Ereignis eines Strahlengangs das Strahlgewicht (z.\,B. eine
    reflektierende oder absorbierende letzte Fläche), rechnet das Plugin mit aufgezeichneten
    Strahlwegen, deren Speicherbedarf mit der Strahlenzahl wächst, und begrenzt die Rechnung auf
    80 Pupillenringe (etwa 19\,400 Strahlen je Feld und Wellenlänge). Das Log weist darauf hin.
\end{itemize}

\section{Optimieren mit Live-Aktualisierung}\label{sec:rt-optimieren}
\index{Optimierung!Raytracer}\index{Gütefunktion}

\gui{Optimise} startet den Optimierer von Raytatouille mit den Variablen und der Gütefunktion
(\emph{merit function}) des Systems. Was optimiert wird, steht in der Systemdatei: Werte mit
\code{\dq{}variable\dq{}: true} sind Variablen; der Abschnitt \param{optimization} nennt Zielgrößen
(\param{operands}, z.\,B. eine Brennweite mit Ziel und Gewicht) und weitere Beiträge zur
Gütefunktion (\param{generators}, z.\,B. der RMS-Fleck). Im Panel gibt es dafür keine
Parameter.

\begin{schritte}
  \item Laden Sie das System.
  \item Öffnen Sie das Dock \gui{Plot}, um den Verlauf zu verfolgen.
  \item Klicken Sie auf \gui{Optimise}.
  \item Beobachten Sie die Ansichten: Nach jeder angenommenen Iteration werden die geänderten
    Flächen und Linsenränder neu gezeichnet. Im Dock \gui{Plot} wächst die Kurve
    \enquote{\texttt{Merit function}} ($\log_{10}\varphi$ über der Iteration); die Statuszeile
    zeigt \enquote{\texttt{iteration }$k$\texttt{: }$\varphi$\texttt{ = …}}.
  \item Warten Sie das Ende ab -- oder brechen Sie mit \keys{Esc} ab. Der Abbruch endet mit dem
    letzten angenommenen Stand.
\end{schritte}

Am Ende meldet die Statuszeile den Status des Optimierers, die Zahl der Iterationen und den
Endwert der Gütefunktion, in der Form \enquote{\textit{Status} \texttt{after }\textit{n}\texttt{ iterations, }$\varphi$\texttt{ = …}}. Der Lauf heißt im Dock \gui{Results} \enquote{\texttt{Optimisation}}.

Das optimierte System ist ein neuer \emph{Systemstand}. Alle folgenden Aktionen
(\gui{Trace rays}, \gui{Camera image}, \gui{Save system}) verwenden ihn. Die ganze Optimierung ist
ein Rückgängig-Schritt: \menu{Edit,Undo} stellt Zeichnung \emph{und} System von vorher wieder her
(\autoref{sec:rt-undo}).

\begin{warnung}
Der optimierte Stand lebt im Plugin-Prozess. Er wird nicht in die Systemdatei geschrieben, bis
Sie \gui{Save system} wählen. Wird das Plugin neu gestartet oder das Dokument später wieder
geöffnet, lädt das Plugin bei der nächsten Aktion die Systemdatei von der Platte -- also den
\emph{nicht} optimierten Stand -- und zeichnet sie neu. Speichern Sie ein Optimierungsergebnis
deshalb immer mit \gui{Save system}, oder arbeiten Sie mit einem Node-Workflow, der das Ergebnis
im Dokument hält (\autoref{sec:rt-workflow}).
\end{warnung}

\section{Das System speichern}\label{sec:rt-speichern}

\gui{Save system} schreibt das aktuelle System als \datei{.rtt.json}:
\begin{itemize}
  \item an den Ort in \gui{Save as} (Gruppe \gui{Optimise}; der Dateidialog fragt nach einem
    Speicherort für \enquote{Raytatouille system (*.rtt.json)}),
  \item oder, wenn \gui{Save as} leer ist, neben die geladene Datei. Ihr Name endet dann
    auf \datei{.optimised.rtt.json} statt auf \datei{.rtt.json}: Aus der Datei
    \datei{lens.rtt.json} wird \datei{lens.optimised.rtt.json}.
\end{itemize}
Die Statuszeile meldet \enquote{\texttt{saved …}} mit dem vollständigen Pfad.

Enthält das Dokument einen Node-Workflow mit dem System, exportiert \gui{Save system} das System
des Workflows -- einschließlich aller dort gemachten Änderungen.

\section{Rückgängig machen}\label{sec:rt-undo}
\index{Rückgängig!Raytracer}

Jede Aktion des Plugins ist in \produkt{} \textbf{ein} Rückgängig-Schritt (\keys{Strg+Z}),
auch wenn sie abgebrochen wird oder fehlschlägt. Damit auch das optische System -- nicht nur
die Zeichnung -- dem Rückgängig-Machen folgt, nummeriert das Plugin jeden Systemstand und
speichert die Nummer im Attribut \param{rtt.revision} der ersten Fläche. Macht
\menu{Edit,Undo} eine Optimierung rückgängig, stellt die Anwendung auch dieses Attribut wieder
her, und die nächste Aktion rechnet wieder mit dem System \emph{vor} der Optimierung.
\menu{Edit,Redo} stellt den optimierten Stand wieder her.

Mit einem Node-Workflow wird statt der Revisionsnummer die Kennung des Workflow-Stands
(\param{rtt.workflow}) verwendet; eine Optimierung, die ihr Ergebnis in den Workflow schreibt,
wird zusammen mit ihrer Zeichnung rückgängig gemacht.

\section{Der Raytracer im Node-Workflow}\label{sec:rt-workflow}
\index{Node-Workflow!Raytracer}

Der Node-Workflow (\autoref{kap:workflow}) stellt ein optisches System als Blöcke und
Verbindungen dar und kann darin Strahlen und Felder rechnen, Gitter anderer Plugins einbinden und
Kamerabilder in Anzeigeblöcken zeigen. Der Raytracer spielt dabei drei Rollen:

\begin{enumerate}
  \item Er \textbf{erzeugt Workflows} aus einem System (\gui{Workflow from system}) oder neu
    (\gui{New Raytatouille system}).
  \item Er \textbf{zeichnet} das System, das der Workflow beschreibt, und hält die Zeichnung
    aktuell (\gui{Follow the workflow}, \gui{Update drawing}).
  \item Er ist \textbf{Block-Anbieter}: Die Raytatouille-Blöcke (\code{rtt.*}) werden in seinem
    Prozess berechnet.
\end{enumerate}

\subsection{Workflow aus einem System}

\begin{schritte}
  \item Laden Sie ein System mit \gui{Load system}.
  \item Wählen Sie \gui{Workflow from system}.
  \item Das Fenster \gui{Node workflow} öffnet sich mit den Blöcken: einer Quelle
    (\block{rtt.source}), Abschnitten (\block{rtt.section}, z.\,B. ein Objektiv),
    Spiegeln, Strahlteilern, Detektoren -- und je Strahlengang des Systems einem
    \emph{Zweig} (\emph{branch}).
\end{schritte}

Die Statuszeile meldet die Zahl der Blöcke, Verbindungen und Zweige. Gab es schon einen Workflow
desselben Systems, übernimmt \gui{Workflow from system} Ihre Änderungen daran: Positionen der
Blöcke, stummgeschaltete Blöcke, Parameter, ersetzte Blöcke (mit ihren Anschlüssen),
hinzugefügte Blöcke und Verbindungen, Darstellungen (Transport) der Verbindungen und eigene
Zweige. Die Meldung nennt, was übernommen wurde (\enquote{\texttt{kept …}}).

\subsection{Ein neues System ohne Datei}

\gui{New Raytatouille system} legt ein System direkt im Workflow an: eine Quelle mit
Standardeinstellungen (mm/\si{\micro\metre}, d-Linie, Eintrittspupille \SI{4}{\milli\metre},
ein Feldpunkt auf der Achse -- änderbar im \gui{Parameters…} der Quelle), eine Blende, einen
Detektor und den Zweig \code{main}. Dazwischen fügen Sie Teile aus der Bibliothek ein
(\gui{Add block}: \emph{Raytatouille mirror}, \emph{beam splitter}, \emph{lens},
\emph{detector}), verbinden ihre Anschlüsse und zeichnen den Zweig neu auf. Ein solches System
wird gezeichnet, gerechnet und gespeichert wie ein importiertes. \menu{Edit,Undo} bringt den
vorherigen Workflow zurück.

\subsection{Der Workflow ist maßgeblich}

Sobald der Workflow das System enthält, ist er für den Raytracer maßgeblich:
\begin{itemize}
  \item \gui{Trace rays}, \gui{Camera image} und \gui{Optimise} rechnen mit dem System, das der
    Workflow beschreibt; die Glaskataloge reisen mit.
  \item \gui{Parameters…} eines Raytatouille-Blocks zeigt je Element einen Reiter mit den
    exakten Daten (Radien, Kegelschnittkonstanten, Aperturen, Lagen der Flächen, Gläser,
    Wechselwirkungen), Längen in mm, Winkel in Grad; die Eigenschaft \enquote{variable} eines
    Werts ist ein Kontrollkästchen. Nur die Werte, die Sie ändern, werden zurückgeschrieben.
  \item Kurz nach jeder Änderung im Workflow zeichnet das Plugin das System neu
    (\gui{Follow the workflow} -- die Anwendung startet diese Aktion selbst, wenn das Plugin
    gerade nichts rechnet). Ein halb eingegebener Wert wartet einfach auf die nächste Änderung;
    Probleme stehen im Log (\enquote{\texttt{drawing not updated: …}}).
  \item \gui{Update drawing} zeichnet sofort. Vor jeder Rechnung wird ohnehin neu gezeichnet,
    wenn sich der Workflow geändert hat. Ohne Workflow meldet die Aktion \enquote{\texttt{No
    node workflow with a system – Workflow from system makes one.}}
  \item \gui{Optimise} schreibt sein Ergebnis in die Blöcke des Workflows.
  \item \gui{Save system} exportiert das System des Workflows als \datei{.rtt.json}.
\end{itemize}

Ohne Workflow arbeitet das Plugin wie in den vorigen Abschnitten mit der geladenen Datei.

\begin{hinweis}
Ein Workflow aus einer älteren Version, dessen Blöcke die Systemdaten noch nicht enthalten, gilt
nicht als maßgeblich; das Plugin rechnet dann mit der Systemdatei. Führen Sie einmal
\gui{Workflow from system} aus, um die Daten in den Workflow zu übernehmen.
\end{hinweis}

\subsection{Abschnitt in eine Gruppe zerlegen}

Ein Block \block{rtt.section} fasst mehrere Elemente einer Baugruppe zusammen, z.\,B. ein ganzes
Objektiv. Im Kontextmenü eines solchen Blocks bietet der Raytracer \gui{Section $\to$ group}
an: Der Abschnitt wird eine Gruppe mit einem Block je Element (Linsengruppe, Blende, \dots), jeder
mit seinen Daten und seinem Teil jedes Durchgangs. Das Licht bleibt dasselbe; die Aktion ist ein
Rückgängig-Schritt. Ein Doppelklick öffnet die Gruppe. Danach können Sie z.\,B. eine Variante der
Gruppe anlegen und darin eine Linsengruppe durch eine ideale Linse ersetzen (siehe
\autoref{kap:workflow}).

\subsection{Raytatouille-Blöcke}

Der Raytracer liefert dem Workflow die Blocktypen \code{rtt.*}: Quelle, Abschnitt, Spiegel,
Strahlteiler, Detektor und die neuen Teile \code{rtt.part.*} (Spiegel, Strahlteiler, Linse,
Blende, Detektor). Die Workflow-Maschine (Plugin \plugin{Workflow}) fragt den Raytracer nach
diesen Typen und lässt sie in seinem Prozess berechnen; jeder Lauf wird einmal mit dem ganzen
Graphen vorbereitet (System aus den Daten der Blöcke, Glaskataloge der Quelle). Die Tabelle aller
Blocktypen mit Anschlüssen und Parametern steht in \autoref{sec:ref-rt-bloecke}.

\subsection{Felder statt Strahlen (Darstellung H)}\label{sec:rt-felder}
\index{Fieldtracing}\index{Darstellung H}

Jede Verbindung des Workflows trägt Licht in einer \emph{Darstellung}: \textbf{R} für Strahlen,
\textbf{H} für Felder. Felder rechnet das Paket \code{fieldtracing}; ist es im Python der Plugins
installiert, kann der Planer des Workflows Verbindungen als H planen bzw. eine auf H festgelegte
Verbindung ausführen. Ohne \code{fieldtracing} stehen nur Strahlen zur Verfügung.

Felder braucht man dort, wo Strahlen nicht genügen, vor allem in der Nähe eines Fokus. Das Panel
des Plugins \plugin{Workflow} enthält dazu \gui{Field accuracy $\boldsymbol{\varepsilon}$ (H links)} (die erlaubte
Abweichung einer Feldausbreitung relativ zur Spitzenamplitude, Standard 0{,}02) und \gui{Wave
optics near a focus (automatic)} (Standard: an): Eine Kamera an einer Verbindung auf
\enquote{auto} prüft, ob sich die Strahlen in ihrer Nähe kreuzen oder fokussieren, und rechnet
dann ab der letzten Fläche mit Feldern. Nach dem Lauf zeigt die Verbindung, was verwendet wurde,
z.\,B. \textbf{R$\to$H}. Einzelheiten stehen in \autoref{kap:workflow}.

\section{Typische Arbeitsabläufe}\label{sec:rt-ablaeufe}

\subsection{Eine Einzellinse optimieren}\label{sec:rt-ablauf-singlet}

Das mitgelieferte Beispiel zeigt den ganzen Ablauf in wenigen Minuten.

\begin{schritte}
  \item \menu{File,New} -- ein leeres Dokument.
  \item Panel \gui{Raytracer}: \gui{System file} $\to$
    \datei{examples/singlet\_optim.rtt.json} aus dem Ordner des Plugins.
  \item \gui{Load system}. Es erscheinen Blende (Ring), Linse \code{L1} mit Rand und die
    Bildfläche \code{IMG}. Der Modellbaum enthält \enquote{Raytatouille: Singlet to optimize: …}.
  \item \gui{Trace rays} mit den Standardwerten: 127 Strahlen. Der Fokus liegt noch neben der
    Bildfläche.
  \item Dock \gui{Plot} öffnen, \gui{Optimise}. Die Linsenradien und der Abstand der Bildfläche
    ändern sich live; die Gütefunktion fällt.
  \item Noch einmal \gui{Trace rays}: Die Strahlen treffen sich jetzt auf der Bildfläche.
  \item \gui{Camera image}: das Fleckbild auf \code{IMG} im Fenster \gui{Image}. Schalten Sie
    \gui{log} ein, um den Rand des Flecks zu sehen.
  \item \gui{Save system}: Das Ergebnis liegt als \datei{singlet\_optim.optimised.rtt.json} neben
    der Ausgangsdatei.
  \item Gefällt das Ergebnis nicht: \keys{Strg+Z} macht die Optimierung rückgängig, Zeichnung und
    System sind wieder im Ausgangszustand.
\end{schritte}

\subsection{Ein Objektiv mit Katalogglas analysieren}

\begin{schritte}
  \item Legen Sie die AGF-Kataloge in den Ordner der Systemdatei oder in seinen Unterordner
    \datei{catalogs} -- oder fügen Sie sie unter \gui{Glass catalogues} hinzu.
  \item \gui{Load system}; prüfen Sie im \gui{Plugin log}, welche Kataloge geladen wurden.
  \item \gui{Trace rays} mit \gui{Field} = \code{0}, \gui{Wavelengths} = \code{all},
    \gui{Colour by} = \code{field}: Die Bündel der Feldpunkte sind unterscheidbar.
  \item Für die Farbfehler: \gui{Colour by} = \code{wavelength} und \gui{Field} = \code{1}
    (nur der Achspunkt).
  \item \gui{Camera image} mit \gui{Field} = \code{2} zeigt das Fleckbild nur des zweiten
    Feldpunkts; mit \gui{Wavelengths} = \code{reference} nur bei der Referenzwellenlänge.
  \item Blenden Sie im Modellbaum vordere Linsen aus, um die Strahlen im Inneren zu sehen.
\end{schritte}

\subsection{Ein Interferometer mit mehreren Strahlengängen}

\begin{schritte}
  \item Laden Sie das System. Die Statuszeile von \gui{Trace rays} nennt jeden Strahlengang mit
    seinen Strahlen.
  \item \gui{Colour by} = \code{path}: Jeder Arm hat seine Farbe; im Dock \gui{Results} lässt
    sich jeder Arm einzeln ein- und ausblenden.
  \item \gui{Paths} = \textit{Name eines Arms}: nur dieser Arm.
  \item \gui{Camera image}: Arme, die auf derselben Fläche enden, werden addiert (inkohärent).
    Interferenz zwischen den Armen berechnet der Node-Workflow (\autoref{kap:workflow}).
  \item Meldet das Plugin, dass es nicht in einen Strahlengang zielen kann, prüfen Sie, ob die
    Aperturblende die erste Fläche dieses Strahlengangs ist (\autoref{sec:rt-gefaltet}).
\end{schritte}

\subsection{Ein System im Workflow aufbauen und weitergeben}

\begin{schritte}
  \item \gui{New Raytatouille system} -- oder ein vorhandenes System laden und \gui{Workflow
    from system}.
  \item Im Fenster \gui{Node workflow} Teile einfügen und verbinden; mit \gui{Parameters…} ihre
    Lage und Daten eingeben. Die CAD-Ansichten zeigen nach jeder Änderung das neue System.
  \item \gui{Check workflow} (Plugin \plugin{Workflow}): Der Workflow prüft das System und
    speichert die Glaskataloge im Dokument.
  \item \gui{Run workflow} rechnet die Anzeigeblöcke; \gui{Trace rays} und \gui{Camera image}
    im Raytracer-Panel rechnen mit demselben System.
  \item Dokument speichern (\keys{Strg+S}). Die \datei{.rttf}-Datei enthält das System, die
    Kataloge und die letzten Ergebnisse der Anzeigeblöcke; sie rechnet auch auf einem anderen
    Rechner ohne die AGF-Dateien.
  \item Wer das System in Raytatouille selbst weiterbearbeiten will, bekommt es mit
    \gui{Save system} als \datei{.rtt.json}.
\end{schritte}

\section{Fehlerbilder und Abhilfe}\label{sec:rt-fehler}

\begin{referenz}{@{}>{\raggedright\arraybackslash}p{0.38\linewidth} >{\raggedright\arraybackslash}X@{}}
\toprule
Meldung oder Beobachtung & Ursache und Abhilfe \\
\midrule\endhead
\texttt{Choose the optical system first: …} & Kein \gui{System file} gewählt. Datei wählen, dann \gui{Load system}. \\
\texttt{The system file does not exist: …} & Datei verschoben oder gelöscht (z.\,B. Dokument von einem anderen Rechner). Pfad neu wählen. \\
\textit{Datei} \texttt{cannot be used: …} & Raytatouille kann das System nicht übersetzen, häufig wegen eines fehlenden Glases. Kataloge prüfen (\autoref{sec:rt-glaskataloge}); sonst die Fehlermeldung von Raytatouille lesen. \\
\texttt{The Python package 'raytatouille' is not installed …} & Paket im Python der Plugins installieren (\gui{Check environment}). \\
\texttt{The irradiance needs a newer Raytatouille …} & Raytatouille aktualisieren. \\
\texttt{Unknown path …} & Name unter \gui{Paths} falsch geschrieben; die Meldung nennt die vorhandenen. \\
\texttt{Raytatouille cannot aim rays into this path …} & Gefalteter oder gekippter Strahlengang, dessen erste Fläche nicht die Aperturblende ist (\autoref{sec:rt-gefaltet}). \\
\texttt{No camera image: none of the paths ends on a drawn plane surface.} & Die letzte Fläche ist gekrümmt oder nicht gezeichnet. Eine ebene Bildfläche ins System aufnehmen. \\
Kamerabild fast leer, viele Strahlen \enquote{outside the grid} & Die Bildfläche ist kleiner als der Fleck. Apertur der Bildfläche im System vergrößern. \\
Fleck im Kamerabild kleiner als ein Pixel & \gui{Pixels (long side)} erhöhen oder die Bildfläche mit kleinerer Apertur definieren. \\
Kamerabild verrauscht & \gui{Rays per field and wavelength} erhöhen. \\
Nach dem Öffnen eines Dokuments ist das Optimierungsergebnis weg & Der optimierte Stand war nur im Plugin-Prozess; das Plugin hat die Datei neu geladen. Vor dem Schließen \gui{Save system} wählen oder mit Workflow arbeiten. \\
\texttt{The node workflow needs the Python package rtt\_fusion …} & \code{rtt-fusion} im Python der Plugins installieren. \\
\texttt{A new system needs a newer RTT-Fusion …}, \texttt{Section $\to$ group needs a newer RTT-Fusion …} & RTT-Fusion aktualisieren. \\
\texttt{The system of the node workflow is not valid: …} & Ein Wert im Workflow ergibt kein gültiges System. Letzte Änderung prüfen oder mit \keys{Strg+Z} zurücknehmen. \\
\texttt{Section $\to$ group needs a node workflow made from a system …} & Zuerst \gui{Workflow from system}. \\
Ein Element ist im CAD als inaktiv markiert, wird aber trotzdem gerechnet & Das Attribut \param{sim.inactive} wird vom Raytracer noch nicht ausgewertet; das Log nennt solche Elemente (\enquote{\texttt{marked inactive in CAD but still traced …}}). \\
\bottomrule
\end{referenz}

\section{Grenzen}\label{sec:rt-grenzen}

\begin{itemize}
  \item Elemente lassen sich noch nicht für die Rechnung abschalten. Objekte mit dem
    Attribut \param{sim.inactive} (im Modellbaum grau dargestellt) werden trotzdem gerechnet
    und nur im Log genannt.
  \item Optische Parameter (Radien, Dicken) lassen sich nicht über die CAD-Geometrie bearbeiten,
    sondern nur in der Systemdatei oder im Node-Workflow.
  \item Die Bestrahlungsstärke wird nur auf ebenen letzten Flächen abgebildet.
  \item Die Ergebnisse der Panel-Aktionen (Strahlen, Kamerabilder, Optimierungskurve) werden
    nicht im Dokument gespeichert; gespeichert werden die Ergebnisse der Anzeigeblöcke des
    Node-Workflows.
\end{itemize}
